% Table: Access to and take-up of A-level economics by school type (Chapter 4 appendix)
\begin{table}[htbp]
\centering
\caption{Access to and take-up of A-level economics by school type}
\label{tab:school_type}
\begin{tabular}{l ccccc}
\toprule
                              & \% offer   & \% offer   & Take-up    & Take-up    & Conversion \\
                              & econ       & maths      & $|$ offer  & $|$ offer  & to econ    \\
                              &            &            & (raw)      & (+ att.)   & degree     \\
\midrule
\multicolumn{6}{c}{\textbf{Panel A: Boys}} \\
\midrule
Mixed                         & 67.8         &    99.1      & 14.1         & XX         & 33.4         \\
Single-sex boys               & 86.9         & 99.5         & 20.0         & XX         & 38.3         \\
\midrule
\multicolumn{6}{c}{\textbf{Panel B: Girls}} \\
\midrule
Mixed                         & 66.4         & 98.6         & 4.8         & XX         & 30.2         \\
Single-sex girls              & 63.1         & 99.3         & 10.6         & XX         & 33.7         \\
\bottomrule
\end{tabular}
\begin{minipage}{\textwidth}
\smallskip
\footnotesize
\textit{Notes:} The table reports the share of schools offering A-level economics and maths, the take-up rate conditional on the school offering economics (raw and controlling for GCSE attainment), and the conversion rate from A-level economics to an economics degree. Sample: 2011 GCSE cohort.
\end{minipage}
\end{table}
